%!TEX root =  ../article.tex

%%%%%%%%%%%%%%%%%%%%%%%%%%%%%%%%%%%%%%%%%%%%%%%%%%%%%%%%%%%%%%%%%%%%%%%%%%%%%%
\section{Related Work} \label{sec:relworks}
%%%%%%%%%%%%%%%%%%%%%%%%%%%%%%%%%%%%%%%%%%%%%%%%%%%%%%%%%%%%%%%%%%%%%%%%%%%%%%
Taskgrind relies on Valgrind to instrument memory accesses and attach them to segments, it eventually implements a determinacy race analysis based on these information.
We present a brief review of the literature on programs instrumentation framework and determinacy races analysis.

\paragraph{Instrumentation Frameworks}
Instrumenting programmes to analyse them is a much-studied subject in the literature, and several approaches had been explored.
Pin~\cite{pin} and Valgrind~\cite{valgrind} rely on Dynamic Binary Translation (DBT)~\cite{dbt}: the binary program is translated to an intermediate representation, instrumented, and its execution is emulated keeping track of a virtual machine state.
The main drawback is the cost of translating and emulating, that lead to important slowdown when executing.
It motivates the design of instrumentation frameworks such as DynamoRIO~\cite{dynamorio}, Dyninstr~\cite{dyninstr} or Instrew~\cite{llvm-dbi}, which is capable of instrumenting while executing natively.
All these frameworks presents different capabilities, in the end, we choose Valgrind for its well-established capabilities (memory allocator overloading, memory accesses instrumentation, debug information reading...) and for its active community.

\paragraph{Determinacy Race Analyzer}
In 1997, Nondeterminator~\cite{nondeterminator} was designed as an entire toolsuite (including compile-time instrumentation and run-time analysis) to detect accurately and provably determinacy races of Cilk programs.
The core of its analyzer relies on a low complexity algorithm (SP-Bags) which assumes the program serial execution correct - also known as the \emph{serial elision}.
Taskgrind has no such assumption.

Helgrind$^{+}$ (2007-2009) is a Valgrind tool for data race detection.
It supports pthread synchronizations primitives, including mutexes and condition variables.
Taskgrind does not support these primitives but targets task-based synchronizations such as barriers or dependencies.
Adding full pthread-support is left for future work.

ThreadSanitizer~\cite{thread-sanitizer-valgrind} (2009) was originally designed as a Valgrind tool, support pthread-based synchronizations.
Its core moved to the LLVM compiler~\cite{tsan-llvm} to reduce execution overheads by instrumenting at compile-time and executing natively.

R. Raghavan et al.~\cite{x10-determinacy-race} (2010) generalized Nondeterminator approach for async/await programming models, and provided a tool for X10 programs, still under the serial elision assumption.

Archer~\cite{archer} was introduced in 2016 as a ThreadSanitizer extension to support OpenMP semantics.
However, it is a thread-centric analyzer that may be a source of false negatives running under task-related synchronizations which is a source of motivations for TaskSanitizer~\cite{Matar18}, ROMP~\cite{ROMP} and Taskgrind.

OmpSs-2 has a toolchain to verify the parallelization of programs~\cite{ompss2-correctness}.
It checks five errors (E1, E2, E3, E4, E5) "\emph{using local task analysis to deal with each task separately}" - as opposed to Taskgrind which looks at the program entirely through its segment graph.
If OmpSs-2 were to be supported by Taskgrind, errors E3, E4 and E5 would fallback to a generic determinacy-error.
On the other hand, OmpSs-2 toolchain seems to provide more insights about the error, such as synchronizations mechanism suggestions.
Extending Taskgrind to report programming model-specific suggestion is left as future work, and could assist programmers in parallelizing their applications.

	   
